\chapter{Master Fields}
\label{chap:master-field}
\chapterepigraph{\jp{三十輻，共一轂，當其無，有車之用。}}{Laozi, \textit{Tao Te Ching}, Chapter 11}

\section{From geometry to a single field}

The tortoise coordinate maps the exterior region to the real line of a scattering problem. The ten components of a \term{metric perturbation} $h_{\mu\nu}$, however, cannot simply be scattered as they stand. One must separate components removable by coordinate transformations, components fixed by \term{constraint equations}, and components that propagate. A variable---or a small set of variables---that packages the resulting physical content is called a \term{master field}.

We first see the mechanism for a massless scalar field. Take the action to be
\begin{equation}
I[\Phi]
=
-\frac{1}{2}
\int \rmd^4x\sqrt{-g}\,
g^{\mu\nu}\partial_\mu\Phi\partial_\nu\Phi
+\int \rmd^4x\sqrt{-g}\,J\Phi,
\label{eq:scalar-action}
\end{equation}
where $J$ is an external source. Expanding in spherical harmonics and setting $\Phi=r^{-1}\Psi_{\ell m}Y_{\ell m}$, the action after angular integration is, up to boundary terms,
\begin{equation}
I_{\ell m}
=
\frac{1}{2}
\int \rmd t\rmd \rstar
\left[
(\partial_t\Psi_{\ell m})^2
-(\partial_{\rstar}\Psi_{\ell m})^2
-V_\ell\Psi_{\ell m}^2
+2S_{\ell m}\Psi_{\ell m}
\right].
\label{eq:reduced-master-action}
\end{equation}
Variation gives
\begin{equation}
\left(
\partial_t^2-\partial_{\rstar}^2+V_\ell
\right)\Psi_{\ell m}
=
S_{\ell m}.
\label{eq:master-action-eom}
\end{equation}
We retain not only the one-dimensional equation but also the action because the same field will later be quantized.

\section{The minimum needed for gravitational perturbations}

Let $g_{\mu\nu}$ be the background metric and perturb it as
\begin{equation}
g_{\mu\nu}
\longrightarrow
g_{\mu\nu}+h_{\mu\nu}.
\end{equation}
Under an \term{infinitesimal coordinate transformation} $x^\mu\to x^\mu+\xi^\mu$,
\begin{equation}
h_{\mu\nu}
\longrightarrow
h_{\mu\nu}
-\nabla_\mu\xi_\nu
-\nabla_\nu\xi_\mu.
\label{eq:linear-gauge-transformation}
\end{equation}
Individual components of $h_{\mu\nu}$ therefore cannot by themselves be interpreted as physical response. On a Schwarzschild background, decomposition by \term{parity} on the sphere separates the perturbation into \term{axial} and \term{polar} sectors. After the constraints are solved, one \term{gauge-invariant} master field remains in each sector for every $\ell\geq2$ \cite{Regge:1957td,Zerilli:1970wzz,Moncrief:1974am}.

Under inversion on the sphere, $(\vartheta,\varphi)\mapsto(\pi-\vartheta,\varphi+\pi)$, the \term{polar} (even-parity) components transform as $(-1)^\ell$, while the \term{axial} (odd-parity) components transform as $(-1)^{\ell+1}$. One \term{radiative degree of freedom} appears in each sector only for $\ell\geq2$. The $\ell=0$ polar perturbation represents an infinitesimal change of mass plus gauge, while the $\ell=1$ axial perturbation represents an infinitesimal change of angular momentum---that is, the linearized Kerr direction---plus gauge. The $\ell=1$ polar vacuum perturbation is a gauge mode corresponding to a displacement of the center of mass. In the presence of matter sources, constrained \term{monopole and dipole} fields remain, but they are not \term{gravitational radiation} propagating to infinity \cite{Martel:2005ir}. Thus the restriction $\ell\geq2$ in the \term{Regge--Wheeler--Zerilli wave equations} below reflects the existence of radiative degrees of freedom, not a convenience of the potentials.

For the axial \term{Regge--Wheeler field} $\Psi_{\mathrm{RW}}$,
\begin{equation}
\left(
\partial_t^2-\partial_{\rstar}^2+V_{\mathrm{RW}}
\right)\Psi_{\mathrm{RW}}
=
S_{\mathrm{RW}},
\label{eq:rw-master}
\end{equation}
where
\begin{equation}
V_{\mathrm{RW}}(r)
=
f(r)
\left[
\frac{\ell(\ell+1)}{r^2}
-\frac{6M}{r^3}
\right],
\qquad
f(r)=1-\frac{2M}{r}.
\label{eq:rw-potential}
\end{equation}
The polar \term{Zerilli field} satisfies an equation of the same form. Defining
\begin{equation}
\lambda
=
\frac{(\ell-1)(\ell+2)}{2},
\end{equation}
its potential is
\begin{align}
V_{\mathrm Z}(r)
&=
\frac{2f(r)}{r^3(\lambda r+3M)^2}
\Bigl[
\lambda^2(\lambda+1)r^3
+3\lambda^2Mr^2
\nonumber\\
&\hspace{9em}
+9\lambda M^2r
+9M^3
\Bigr].
\label{eq:zerilli-potential-minimum}
\end{align}
The two potentials are different, but they are related by an appropriate first-order differential transformation and have the same quasinormal-mode frequencies, to be defined in Chapter~\ref{chap:qnm} \cite{Chandrasekhar:book,Chandrasekhar:1975nkd}.

\begin{principlebox}
A master field and its effective potential are representations and are not unique. What should be compared across representations are the radiation conditions, the real-frequency \term{scattering matrix}, and the map from physical sources to observables.
\end{principlebox}

\section{Sources and observables}

Vacuum equations alone do not constitute a response theory. If the matter \term{stress tensor} is $T_{\mu\nu}$, spherical-harmonic components of the \term{linearized Einstein equations} determine $S_{\mathrm{RW}}$ or $S_{\mathrm Z}$. Even for the same physical matter source, changing the definition of the master field changes the form of $S$. The master-field solution itself must therefore not be identified directly with an observable.

Let $O$ denote an \term{observation operation}, such as a far-zone radiative waveform, a flux crossing the horizon, or a local curvature perturbation. Schematically, the relation from source to observable is
\begin{equation}
T_{\mu\nu}
\longrightarrow
S_{\ell m}
\longrightarrow
\Psi_{\ell m}
\longrightarrow
\delta O.
\label{eq:source-master-observable-chain}
\end{equation}
The first arrow represents the constraints and projection of the linearized Einstein equations, while the final arrow represents reconstruction of a waveform or flux. In the next chapter, the central arrow will be defined as a solution operator responding to an external source.

\paragraph{One complete chain.}
With the Martel--Poisson normalization, the gravitational wave at future \term{null infinity} can be reconstructed from the polar field $\Psi^{\mathrm{even}}_{\ell m}$ and axial field $\Psi^{\mathrm{odd}}_{\ell m}$ as
\begin{equation}
h_+-\rmi h_\times
=
\frac{1}{2r}
\sum_{\ell m}
\sqrt{\frac{(\ell+2)!}{(\ell-2)!}}
\left(
\Psi^{\mathrm{even}}_{\ell m}
+\rmi\Psi^{\mathrm{odd}}_{\ell m}
\right)
{}_{-2}Y_{\ell m}
+O(r^{-2}).
\label{eq:master-to-strain}
\end{equation}
With $u=t-\rstar$, the radiated energy is
\begin{equation}
\frac{\rmd E}{\rmd u}
=
\frac{1}{64\pi}
\sum_{\ell m}
\frac{(\ell+2)!}{(\ell-2)!}
\left(
\left|\partial_u\Psi^{\mathrm{even}}_{\ell m}\right|^2
+
\left|\partial_u\Psi^{\mathrm{odd}}_{\ell m}\right|^2
\right).
\label{eq:master-to-energy-flux}
\end{equation}
This normalization follows Ref.~\cite{Martel:2005ir}. Other conventions for the RW field insert a time derivative or integral into the odd-parity field and change numerical coefficients. An observable is therefore obtained only by carrying, in one consistent normalization, the projection of the source $T_{\mu\nu}$, the Green-function calculation of $\Psi_{\ell m}$, and the \term{waveform reconstruction} in Eq.~\eqref{eq:master-to-strain}.

\section{What to retain about rotating spacetimes}

Because \term{Kerr spacetime} is not spherically symmetric, it does not reduce to a single Regge--Wheeler-type real potential. Let the mass be $M$ and the angular momentum be $J=aM$, and define
\begin{align}
\Delta
&=
r^2-2Mr+a^2,
&
r_\pm
&=
M\mathord{\pm}\sqrt{M^2-a^2},
&
\Omega_H
&=
\frac{a}{r_+^2+a^2}.
\label{eq:kerr-horizon-minimum}
\end{align}
Here $r_+$ is the event-horizon radius and $\Omega_H$ is its angular velocity. The Kerr tortoise coordinate is defined by
\begin{equation}
\frac{\rmd\rstar}{\rmd r}
=
\frac{r^2+a^2}{\Delta}.
\label{eq:kerr-tortoise-minimum}
\end{equation}

Using curvature components adapted to \term{Petrov type D}, a spin-$s$ \term{Teukolsky variable} can be separated schematically as
\begin{equation}
\psi_s
=
\rme^{-\rmi\omega t+\rmi m\varphi}
S_{s\ell m}(\vartheta;a\omega)
R_{s\ell m}(r).
\label{eq:teukolsky-separation-minimum}
\end{equation}
\cite{Teukolsky:1972my} Here $S_{s\ell m}$ is a \term{spin-weighted spheroidal harmonic}. Because the \term{angular eigenvalue} depends on $a\omega$, the angular and radial equations must be solved together. Suppressing spin-dependent power-law prefactors, the radiation conditions are
\begin{align}
R_{s\ell m}
&\sim
\rme^{-\rmi(\omega-m\Omega_H)\rstar},
&&r\longrightarrow r_+,
\label{eq:kerr-horizon-bc-minimum}
\\
R_{s\ell m}
&\sim
\rme^{+\rmi\omega\rstar},
&&r\longrightarrow\infty.
\label{eq:kerr-infinity-bc-minimum}
\end{align}
The \term{Killing-energy flux} crossing the horizon is proportional to $\omega-m\Omega_H$. Consequently, for bosons in the range $0<\omega<m\Omega_H$, the horizon flux is negative and the reflected wave is stronger than the incident wave: \term{superradiance} occurs. The same frequency shift reappears in the thermal factor in Chapter~\ref{chap:horizon-thermality}.

The minimum amount of Kerr physics retained in the main text is the \term{horizon generator}, \term{Teukolsky separation}, radiation conditions at both ends, and \term{superradiance}. We do not survey metric reconstruction or the full taxonomy of numerical methods; implementation and verification are deferred to the research problem in Chapter~\ref{chap:qnm}.

\section{Conclusion of this chapter}

Black-hole perturbations reduce, after gauge and constraints have been removed, to an \term{open wave problem} for master fields. The action \eqref{eq:reduced-master-action} is the common starting point for classical response and quantization, while Eq.~\eqref{eq:source-master-observable-chain} separates representation-dependent quantities from observables. In the next chapter we invert the master equation in the presence of an external source.
